\documentclass[10pt,]{letter}
\usepackage{mathpazo}
\usepackage{amssymb,amsmath}
\usepackage{ifxetex,ifluatex}
\usepackage{fixltx2e} % provides \textsubscript
\ifnum 0\ifxetex 1\fi\ifluatex 1\fi=0 % if pdftex
  \usepackage[T1]{fontenc}
  \usepackage[utf8]{inputenc}
\else % if luatex or xelatex
  \ifxetex
    \usepackage{mathspec}
  \else
    \usepackage{fontspec}
  \fi
  \defaultfontfeatures{Ligatures=TeX,Scale=MatchLowercase}
\fi
% use upquote if available, for straight quotes in verbatim environments
\IfFileExists{upquote.sty}{\usepackage{upquote}}{}
% use microtype if available
\IfFileExists{microtype.sty}{%
\usepackage{microtype}
\UseMicrotypeSet[protrusion]{basicmath} % disable protrusion for tt fonts
}{}
\usepackage[left=2cm, right=2cm, top=2cm, bottom=3cm, footskip =
.5cm]{geometry}
\usepackage[unicode=true]{hyperref}
\PassOptionsToPackage{usenames,dvipsnames}{color} % color is loaded by hyperref
\hypersetup{
            pdfauthor={Joseph Caracappa, PhD; Research Fishery Biologist; Ecosystem Dynamics and Assessment Branch; Northeast Fisheries Science Center},
            colorlinks=true,
            linkcolor=Maroon,
            citecolor=Blue,
            urlcolor=blue,
            breaklinks=true}
\urlstyle{same}  % don't use monospace font for urls
% Make links footnotes instead of hotlinks:
\renewcommand{\href}[2]{#2\footnote{\url{#1}}}
\IfFileExists{parskip.sty}{%
\usepackage{parskip}
}{% else
\setlength{\parindent}{0pt}
\setlength{\parskip}{6pt plus 2pt minus 1pt}
}
\setlength{\emergencystretch}{3em}  % prevent overfull lines
\providecommand{\tightlist}{%
  \setlength{\itemsep}{0pt}\setlength{\parskip}{0pt}}
\setcounter{secnumdepth}{0}

% set default figure placement to htbp
\makeatletter
\def\fps@figure{htbp}
\makeatother

\usepackage{graphicx,grffile}
\signature{\vspace*{-8ex}\includegraphics{C:/Users/joseph.caracappa/Documents/GitHub/READ-EDAB-SOE\_reports/utils/caracappa\_signature.pdf}\\Joseph
Caracappa, PhD\\Research Fishery Biologist\\Ecosystem Dynamics and
Assessment Branch\\Northeast Fisheries Science Center}

\date{19 March, 2026}

\usepackage{float}
\address{}

\usepackage{mdframed} % color is loaded by mdframed
\definecolor{greyborder}{RGB}{221,221,221}
\definecolor{greytext}{RGB}{119,119,119}
\newmdenv[rightline=false,bottomline=false,topline=false,linewidth=3pt,linecolor=greyborder,skipabove=\parskip]{blockquote}
\renewenvironment{quote}{\begin{blockquote}\list{}{\rightmargin=0em\leftmargin=0em}%
\item\relax\color{greytext}\ignorespaces}{\unskip\unskip\endlist\end{blockquote}}

\usepackage{wallpaper}
\ThisULPositionWallPaper{1}{0pt}{0pt}{utils/NEFSCletterhead.pdf}


\begin{document}


\begin{letter}{New England Fishery Management Council\\50 Water St
No.~2\\Newburyport, MA 01950}
\opening{To the SSC,}

In this memo we list comments and requests received on the 2019-2025
State of the Ecosystem (SOE) reports, and how we responded to those
requests. We include comments from both Councils because adjustments to
the report were made in response to both. We welcome feedback on whether
this memo is useful and how to improve it for future SOE reporting.

This year, the memo is reorganized into three tables. Table 1 shows all
completed requests, Table 2 shows all pending requests, and Table 3
shows all request removed from tracking. Request priorities and the
decision that items were out of scope for the SOE were discussed with
both Councils' SSCs during a July 2025 prioritization meeting.

The SOE team is also requesting that the Council and SSC consider some
of the following items relating to recent changes to the SOE

\begin{enumerate}
\def\labelenumi{\arabic{enumi}.}
\item
  \textbf{Performance metrics}: For each management objective category
  are the performance metrics clearly articulated and are the metrics
  used appropriate?
\item
  \textbf{Specific objectives and indicators}:

  \begin{itemize}
  \tightlist
  \item
    The framework of the \emph{Stability objective} has been revised for
    both ecosystem and fishery indicators to focus on \emph{Changes From
    Baseline}, \emph{Adaptive Capacity}, and \emph{Volatility}.
  \item
    Does this reframing of stability better address the Council needs?
    Comment on how they could be refined
  \end{itemize}
\item
  The SOE team has integrated \textbf{ocean model forecasts} from MOM6
  into an SOE graphic summary

  \begin{itemize}
  \tightlist
  \item
    Are presenting forecast information as probabilities of deviating
    from historical average a useful format for the Council? Are there
    other reference points or ``events'' that could be used instead of
    historical averages?
  \item
    How should uncertainty be discussed alongside observations?
  \item
    Are there any research questions or Council decisions that forecast
    information could assist?
  \end{itemize}
\end{enumerate}

We welcome comments on the entire SOE report as well as information
included in this memo, and look forward to feedback from the SSC and
Council.

\longindentation=0pt
\closing{Sincerely,}
\encl{State of the Ecosystem 2024: Request Tracking Memo}
\cc{Jon Hare}

\end{letter}

\end{document}
